\documentclass[11pt]{article}
\usepackage[margin=1.2in]{geometry}
\usepackage{amsmath, amssymb}
\usepackage{bm}
\usepackage{booktabs}
\usepackage{parskip}
\usepackage{hyperref}
\usepackage{xcolor}
\usepackage{titlesec}
\usepackage{enumitem}

\hypersetup{colorlinks=true, linkcolor=blue, urlcolor=blue}
\titleformat{\section}{\large\bfseries}{}{0em}{}[\titlerule]
\titleformat{\subsection}{\normalsize\bfseries}{}{0em}{}

\title{\textbf{Jacobian-L1 Influence in Graph Neural Networks}\\[0.4em]
       \large How pairwise node influence is computed in \texttt{influence.py}}
\author{}
\date{}

\begin{document}
\maketitle

\section{Definition}

For a trained $k$-layer GCN, the \textbf{Jacobian-L1 influence} of node $y$ on
focal node $x$ is:

\begin{equation}
  I(x,\, y)
  \;=\;
  \sum_{i=1}^{d_k}\;\sum_{f=1}^{d_0}
  \left|\,
    \frac{\partial\, h_x^{(k)}[i]}{\partial\, h_y^{(0)}[f]}
  \,\right|
  \label{eq:def}
\end{equation}

where $h_x^{(k)} \in \mathbb{R}^{d_k}$ is node $x$'s embedding after $k$
message-passing layers, and $h_y^{(0)} \in \mathbb{R}^{d_0}$ is node $y$'s
raw input feature vector. The double sum collapses the full Jacobian to a
single non-negative scalar capturing the \emph{total sensitivity} of $x$'s
prediction to $y$'s input features.

\section{GCN Layer Notation}

Each GCN layer aggregates from the 1-hop neighbourhood with symmetric
degree normalisation:

\begin{equation}
  h_v^{(\ell)}
  \;=\;
  \sigma\!\left(
    W^{(\ell)}\;
    \sum_{u \,\in\, \hat{\mathcal{N}}(v)}
    \tilde{a}_{vu}\; h_u^{(\ell-1)}
  \right),
  \qquad
  \tilde{a}_{vu}
  \;=\;
  \frac{1}{\sqrt{(\deg v + 1)(\deg u + 1)}}
  \label{eq:gcn}
\end{equation}

where $\hat{\mathcal{N}}(v)$ includes self-loops, $W^{(\ell)}$ is the
layer-$\ell$ weight matrix, and $\sigma$ is the activation function (e.g.\
ReLU). Define the pre-activation at node $v$ and layer $\ell$ as:

\begin{equation}
  \mathbf{z}_v^{(\ell)}
  \;=\;
  W^{(\ell)}\;
  \sum_{u \,\in\, \hat{\mathcal{N}}(v)}
  \tilde{a}_{vu}\; h_u^{(\ell-1)},
  \qquad\text{so}\quad
  h_v^{(\ell)} = \sigma\!\left(\mathbf{z}_v^{(\ell)}\right).
\end{equation}

\section{Concrete Example: 2-Layer GCN on a 3-Node Chain}

\textbf{Graph:}\quad $x \;{=}\; 0 \;-\; y \;{=}\; 1 \;-\; z \;{=}\; 2$

The focal node is $x$, so $k = 2$ and the $k$-hop subgraph contains all
three nodes.

\subsection{Forward Pass}

\paragraph{Layer 1.}
\begin{align}
  h_x^{(1)} &= \sigma\!\bigl(W^{(1)}\bigl(\tilde{a}_{xx}\,h_x^{(0)}
                + \tilde{a}_{xy}\,h_y^{(0)}\bigr)\bigr) \\[4pt]
  h_y^{(1)} &= \sigma\!\bigl(W^{(1)}\bigl(\tilde{a}_{yx}\,h_x^{(0)}
                + \tilde{a}_{yy}\,h_y^{(0)}
                + \tilde{a}_{yz}\,h_z^{(0)}\bigr)\bigr) \\[4pt]
  h_z^{(1)} &= \sigma\!\bigl(W^{(1)}\bigl(\tilde{a}_{zy}\,h_y^{(0)}
                + \tilde{a}_{zz}\,h_z^{(0)}\bigr)\bigr)
\end{align}

\paragraph{Layer 2.}
\begin{equation}
  h_x^{(2)} = \sigma\!\bigl(W^{(2)}\bigl(\tilde{a}_{xx}\,h_x^{(1)}
               + \tilde{a}_{xy}\,h_y^{(1)}\bigr)\bigr)
\end{equation}

Node $z$ never appears directly in $x$'s layer-2 aggregation (it is 2 hops
away), but it \emph{indirectly} influences $x$ through $y$'s hidden state
$h_y^{(1)}$.

\subsection{Derivative for a 1-Hop Neighbour ($y$)}

Applying the chain rule through layer 2 then layer 1:

\begin{equation}
  \frac{\partial\, h_x^{(2)}[i]}{\partial\, h_y^{(0)}[f]}
  \;=\;
  \sum_j
  \underbrace{%
    \frac{\partial\, h_x^{(2)}[i]}{\partial\, h_x^{(1)}[j]}
  }_{\text{layer 2}}
  \cdot
  \underbrace{%
    \frac{\partial\, h_x^{(1)}[j]}{\partial\, h_y^{(0)}[f]}
  }_{\text{layer 1}}
\end{equation}

Expanding each factor using equation~\eqref{eq:gcn}:

\begin{align}
  \frac{\partial\, h_x^{(1)}[j]}{\partial\, h_y^{(0)}[f]}
  &\;=\;
  \sigma'\!\bigl(\mathbf{z}_x^{(1)}[j]\bigr)
  \;\cdot\;
  W^{(1)}[j,f]
  \;\cdot\;
  \tilde{a}_{xy}
  \\[6pt]
  \frac{\partial\, h_x^{(2)}[i]}{\partial\, h_x^{(1)}[j]}
  &\;=\;
  \sigma'\!\bigl(\mathbf{z}_x^{(2)}[i]\bigr)
  \;\cdot\;
  W^{(2)}[i,j]
  \;\cdot\;
  \tilde{a}_{xx}
\end{align}

Combining:

\begin{equation}
  \boxed{
  \frac{\partial\, h_x^{(2)}[i]}{\partial\, h_y^{(0)}[f]}
  \;=\;
  \sum_j\;
  \sigma'\!\bigl(\mathbf{z}_x^{(2)}[i]\bigr)\cdot W^{(2)}[i,j]\cdot\tilde{a}_{xx}
  \;\cdot\;
  \sigma'\!\bigl(\mathbf{z}_x^{(1)}[j]\bigr)\cdot W^{(1)}[j,f]\cdot\tilde{a}_{xy}
  }
\end{equation}

\medskip
The edge weight $\tilde{a}_{xy}$ appears as a multiplicative factor: a
high-degree neighbour $y$ has a small $\tilde{a}_{xy}$, attenuating its
influence on $x$.

\subsection{Derivative for a 2-Hop Neighbour ($z$)}

Node $z$ feeds into $h_y^{(1)}$, which in turn feeds into $h_x^{(2)}$.
The chain rule now routes through $y$'s hidden state:

\begin{equation}
  \frac{\partial\, h_x^{(2)}[i]}{\partial\, h_z^{(0)}[f]}
  \;=\;
  \sum_j
  \underbrace{%
    \frac{\partial\, h_x^{(2)}[i]}{\partial\, h_y^{(1)}[j]}
  }_{\text{layer 2, via }y}
  \cdot
  \underbrace{%
    \frac{\partial\, h_y^{(1)}[j]}{\partial\, h_z^{(0)}[f]}
  }_{\text{layer 1, at }y}
\end{equation}

\begin{align}
  \frac{\partial\, h_y^{(1)}[j]}{\partial\, h_z^{(0)}[f]}
  &\;=\;
  \sigma'\!\bigl(\mathbf{z}_y^{(1)}[j]\bigr)
  \;\cdot\;
  W^{(1)}[j,f]
  \;\cdot\;
  \tilde{a}_{yz}
  \\[6pt]
  \frac{\partial\, h_x^{(2)}[i]}{\partial\, h_y^{(1)}[j]}
  &\;=\;
  \sigma'\!\bigl(\mathbf{z}_x^{(2)}[i]\bigr)
  \;\cdot\;
  W^{(2)}[i,j]
  \;\cdot\;
  \tilde{a}_{xy}
\end{align}

\begin{equation}
  \boxed{
  \frac{\partial\, h_x^{(2)}[i]}{\partial\, h_z^{(0)}[f]}
  \;=\;
  \sum_j\;
  \sigma'\!\bigl(\mathbf{z}_x^{(2)}[i]\bigr)\cdot W^{(2)}[i,j]\cdot\tilde{a}_{xy}
  \;\cdot\;
  \sigma'\!\bigl(\mathbf{z}_y^{(1)}[j]\bigr)\cdot W^{(1)}[j,f]\cdot\tilde{a}_{yz}
  }
\end{equation}

\medskip
\noindent\textbf{Key difference from the 1-hop case:} the activation gate
at layer 1 is now $\sigma'\!\bigl(\mathbf{z}_y^{(1)}\bigr)$ (intermediate
node $y$'s gate) rather than $\sigma'\!\bigl(\mathbf{z}_x^{(1)}\bigr)$.
The gradient flows \emph{backward through $y$'s hidden state} to reach $z$.
Each additional hop introduces another pair of $(\sigma', W, \tilde{a})$
factors, causing the Jacobian to decay exponentially with distance.

\subsection{Comparison at a Glance}

\begin{center}
\renewcommand{\arraystretch}{1.5}
\begin{tabular}{lll}
\toprule
 & \textbf{1-hop neighbour }$y$ & \textbf{2-hop neighbour }$z$ \\
\midrule
Layer-1 gate   & $\sigma'\!\left(\mathbf{z}_x^{(1)}\right)$ & $\sigma'\!\left(\mathbf{z}_y^{(1)}\right)$ \\
Layer-1 weight & $W^{(1)}$ & $W^{(1)}$ \\
Layer-1 edge   & $\tilde{a}_{xy}$ & $\tilde{a}_{yz}$ \\
Layer-2 gate   & $\sigma'\!\left(\mathbf{z}_x^{(2)}\right)$ & $\sigma'\!\left(\mathbf{z}_x^{(2)}\right)$ \\
Layer-2 weight & $W^{(2)}$ & $W^{(2)}$ \\
Layer-2 edge   & $\tilde{a}_{xx}$ & $\tilde{a}_{xy}$ \\
\bottomrule
\end{tabular}
\end{center}

\section{From Jacobian Tensor to Scalar Influence}

\texttt{torch.autograd.functional.jacobian} returns the full tensor
\[
  J \;\in\; \mathbb{R}^{d_k \times |V_\text{sub}| \times d_0},
  \qquad
  J[i,\, n_\text{local},\, f]
  \;=\;
  \frac{\partial\, h_x^{(k)}[i]}{\partial\, h_{n}^{(0)}[f]}
\]
evaluated over the $k$-hop induced subgraph (shape is $|V_\text{sub}|$
instead of $N$ for tractability; outside the subgraph the true influence
is exactly 0). The reduction to per-node scalars:

\begin{equation}
  I(x,\, n)
  \;=\;
  \sum_{i=1}^{d_k}\sum_{f=1}^{d_0}
  \bigl|J[i,\, n_\text{local},\, f]\bigr|
  \;=\;
  \bigl\|J[\,:\,,\, n_\text{local},\, :\,]\bigr\|_1
\end{equation}

In code (from \texttt{influence.py}, line 88):
\begin{verbatim}
    influence_sub = jac.abs().sum(dim=(0, 2))   # shape: [|V_sub|]
\end{verbatim}

The result is scattered back to global node indices, giving a length-$N$
vector where nodes outside the $k$-hop subgraph have $I(x, n) = 0$.

\section{Effective Influence}

The \emph{effective influence} additionally weights each node's score by
the GCN symmetric normalisation factor it would receive if it were a direct
1-hop neighbour of $x$:

\begin{equation}
  I_\text{eff}(x,\, n)
  \;=\;
  I(x,\, n)
  \;\cdot\;
  \tilde{a}_{xn}
  \;=\;
  I(x,\, n)
  \;\cdot\;
  \frac{1}{\sqrt{(\deg x + 1)(\deg n + 1)}}
\end{equation}

This re-weights raw influence by how strongly a direct edge $x \leftarrow n$
would be attenuated by degree, making high-degree neighbours appear less
influential even when their Jacobian norm is large.

\section{Why Autograd, Not Closed Form?}

Computing the Jacobian analytically requires manually unrolling the chain
rule through every layer for every node, accounting for the activation gates
$\sigma'$ evaluated at the actual forward-pass pre-activations. PyTorch's
\texttt{torch.autograd.functional.jacobian} does this automatically by
either:

\begin{itemize}[leftmargin=1.5em]
  \item \textbf{Reverse mode} (\texttt{vectorize=False}): runs one backward
    pass per output dimension $i$ ($d_k$ passes total).
  \item \textbf{Forward mode} (\texttt{vectorize=True}): uses
    \texttt{torch.vmap} over the input dimensions, computing all columns of
    the Jacobian in a single forward sweep — same result, significantly
    faster.
\end{itemize}

\end{document}
